\chapter{تغییر پایه}
\label{ch:change-of-basis}

یک بردار در پایه‌های \enterm{Basis} مختلف مختصات مختلف دارد، و یک تبدیل خطی \enterm{Linear Transformation} هم در پایه‌های مختلف ماتریس‌های مختلف. تبدیلی که در یک پایه پیچیده به نظر می‌رسد، با پایه‌ی مناسب ساده می‌شود. در این فصل می‌بینیم مختصات و ماتریس‌ها وقتی پایه عوض می‌شود چطور تغییر می‌کنند.

\section{مرور: مختصات نسبت به پایه}

یادآوری فصل ۶: اگر $\mathcal{B}=\{\vb_1,\dots,\vb_n\}$ پایه‌ی مرتب \enterm{Ordered Basis} باشد و $\vx=c_1\vb_1+\dots+c_n\vb_n$، بردار مختصات \enterm{Coordinate Vector} $\coord{\vx}{\mathcal{B}}=(c_1,\dots,c_n)$ است.

\begin{example}
مختصات $2+7x+x^2\in P_2$ را نسبت به $\mathcal{B}=\{x+x^2,\ 1+x^2,\ 1+x\}$ بیابید.
\end{example}
\begin{solution}
باید $2+7x+x^2=c_1(x+x^2)+c_2(1+x^2)+c_3(1+x)$. با مقایسه‌ی ضرایب:
$c_2+c_3=2$،\ $c_1+c_3=7$،\ $c_1+c_2=1$؛ پس $c_1=3$، $c_2=-2$، $c_3=4$ و
$\coord{2+7x+x^2}{\mathcal{B}}=(3,-2,4)$. نسبت به پایه‌ی استاندارد \enterm{Standard Basis} $\mathcal{C}=\{1,x,x^2\}$ مختصات $(2,7,1)$ است.
\end{solution}

برای پایه‌ی $\mathcal{B}=\{\vv_1,\dots,\vv_n\}$ از $\R^n$ و $P_{\mathcal{B}}=\begin{bmatrix}\vv_1&\cdots&\vv_n\end{bmatrix}$ داریم $P_{\mathcal{B}}\coord{\vx}{\mathcal{B}}=\vx$؛ چون این همان نوشتن $\vx$ به صورت ترکیب خطی \enterm{Linear Combination} ستون‌هاست. $P_{\mathcal{B}}$ وارون‌پذیر \enterm{Invertible} است (ستون‌هایش مستقل‌اند) و $\coord{\vx}{\mathcal{B}}=P_{\mathcal{B}}^{-1}\vx$.

\section{ماتریس تغییر پایه}

\begin{theorem}[ماتریس تغییر پایه]\label{thm:change-basis}
اگر $\mathcal{B}=\{\vb_1,\dots,\vb_n\}$ و $\mathcal{C}=\{\vc_1,\dots,\vc_n\}$ دو پایه‌ی $V$ باشند، ماتریس یکتای $n\times n$ مثل $P_{\mathcal{C}\leftarrow\mathcal{B}}$ هست که
\begin{keyformula}
\[
\coord{\vx}{\mathcal{C}}=P_{\mathcal{C}\leftarrow\mathcal{B}}\,\coord{\vx}{\mathcal{B}}\qquad(\vx\in V),
\]
\end{keyformula}
و ستون‌های آن مختصات بردارهای پایه‌ی $\mathcal{B}$ نسبت به $\mathcal{C}$ هستند:
$P_{\mathcal{C}\leftarrow\mathcal{B}}=\begin{bmatrix}\coord{\vb_1}{\mathcal{C}}&\cdots&\coord{\vb_n}{\mathcal{C}}\end{bmatrix}$.
به‌علاوه $P_{\mathcal{C}\leftarrow\mathcal{B}}$ وارون‌پذیر است و $(P_{\mathcal{C}\leftarrow\mathcal{B}})^{-1}=P_{\mathcal{B}\leftarrow\mathcal{C}}$.
\end{theorem}
\begin{proof}
اگر $\vx=\sum a_i\vb_i$، طبق خطی بودن نگاشت مختصات (قضیه‌ی \ref{thm:coord-map}):
\[
\coord{\vx}{\mathcal{C}}=\sum a_i\coord{\vb_i}{\mathcal{C}}=\begin{bmatrix}\coord{\vb_1}{\mathcal{C}}&\cdots&\coord{\vb_n}{\mathcal{C}}\end{bmatrix}\begin{bmatrix}a_1\\ \vdots\\ a_n\end{bmatrix},
\]
و $(a_1,\dots,a_n)=\coord{\vx}{\mathcal{B}}$؛ پس ماتریس با این ستون‌ها کار می‌کند. یکتایی: اگر $M\coord{\vx}{\mathcal{B}}=\coord{\vx}{\mathcal{C}}$ برای هر $\vx$، با $\vx=\vb_j$ داریم $\coord{\vb_j}{\mathcal{B}}=\ve_j$ و $M\ve_j$ ستون $j$ام $M$ است؛ پس ستون‌ها همان‌ها هستند. وارون \enterm{Inverse}: از
$\coord{\vx}{\mathcal{B}}=P_{\mathcal{B}\leftarrow\mathcal{C}}\coord{\vx}{\mathcal{C}}=P_{\mathcal{B}\leftarrow\mathcal{C}}P_{\mathcal{C}\leftarrow\mathcal{B}}\coord{\vx}{\mathcal{B}}$
برای هر $\vx$ (و با یکتایی) $P_{\mathcal{B}\leftarrow\mathcal{C}}P_{\mathcal{C}\leftarrow\mathcal{B}}=I$؛ و برعکس‌اش هم همین‌طور.
\end{proof}

\begin{center}
\begin{tikzpicture}[>=Latex, box/.style={draw=primary, fill=primarylight, rounded corners=2mm, minimum width=2.6cm, minimum height=0.9cm, font=\small}]
  \node[box] (V) at (3.2,2.2) {$\vx\in V$};
  \node[box] (B) at (6.4,0) {$\coord{\vx}{\mathcal{B}}\in\R^n$};
  \node[box] (C) at (0,0) {$\coord{\vx}{\mathcal{C}}\in\R^n$};
  \draw[->, thick] (V)--(B) node[midway, above right, font=\scriptsize] {$\coord{\cdot}{\mathcal{B}}$};
  \draw[->, thick] (V)--(C) node[midway, above left, font=\scriptsize] {$\coord{\cdot}{\mathcal{C}}$};
  \draw[->, thick, warncolor] (B)--(C) node[midway, below, font=\small] {$P_{\mathcal{C}\leftarrow\mathcal{B}}$};
\end{tikzpicture}
\end{center}

\begin{note}[در $\R^n$]
چون $\vx=P_{\mathcal{B}}\coord{\vx}{\mathcal{B}}=P_{\mathcal{C}}\coord{\vx}{\mathcal{C}}$، داریم $\coord{\vx}{\mathcal{C}}=P_{\mathcal{C}}^{-1}P_{\mathcal{B}}\coord{\vx}{\mathcal{B}}$؛ یعنی $P_{\mathcal{C}\leftarrow\mathcal{B}}=P_{\mathcal{C}}^{-1}P_{\mathcal{B}}$. در عمل کافی است $[P_{\mathcal{C}}\,|\,P_{\mathcal{B}}]$ را به RREF ببریم: بخش راست می‌شود $P_{\mathcal{C}\leftarrow\mathcal{B}}$.
\end{note}

\begin{example}
برای $\mathcal{B}=\{x+x^2,\ 1+x^2,\ 1+x\}$ و $\mathcal{C}=\{1,x,x^2\}$ در $P_2$، ماتریس‌های $P_{\mathcal{C}\leftarrow\mathcal{B}}$ و $P_{\mathcal{B}\leftarrow\mathcal{C}}$ را بیابید و مختصات $2+7x+x^2$ نسبت به $\mathcal{B}$ را از آن‌ها به دست آورید.
\end{example}
\begin{solution}
ستون‌های $P_{\mathcal{C}\leftarrow\mathcal{B}}$ مختصات پایه‌ی $\mathcal{B}$ نسبت به پایه‌ی استاندارد هستند:
\[
P_{\mathcal{C}\leftarrow\mathcal{B}}=\begin{bmatrix}0&1&1\\1&0&1\\1&1&0\end{bmatrix},\qquad
P_{\mathcal{B}\leftarrow\mathcal{C}}=P_{\mathcal{C}\leftarrow\mathcal{B}}^{-1}=\frac12\begin{bmatrix}-1&1&1\\1&-1&1\\1&1&-1\end{bmatrix}.
\]
(بررسی: اگر $J$ ماتریس تمام‌یک باشد، $(J-I)(\frac12J-I)=I$.) پس
$\coord{2+7x+x^2}{\mathcal{B}}=P_{\mathcal{B}\leftarrow\mathcal{C}}(2,7,1)=(3,-2,4)$، همان جواب بخش قبل.
\end{solution}

\begin{example}
$\vb_1=(1,-3)$، $\vb_2=(-2,4)$، $\vc_1=(-7,9)$، $\vc_2=(-5,7)$ و $\mathcal{B}=\{\vb_1,\vb_2\}$، $\mathcal{C}=\{\vc_1,\vc_2\}$ پایه‌های $\R^2$ هستند. ماتریس تبدیل مختصات از $\mathcal{C}$ به $\mathcal{B}$ و از $\mathcal{B}$ به $\mathcal{C}$ را بیابید.
\end{example}
\begin{solution}
ستون‌های $P_{\mathcal{B}\leftarrow\mathcal{C}}$ مختصات $\vc_1,\vc_2$ نسبت به $\mathcal{B}$ هستند؛ با کاهش سطری
\[
\left[\begin{array}{cc|cc}1&-2&-7&-5\\-3&4&9&7\end{array}\right]\sim
\left[\begin{array}{cc|cc}1&0&5&3\\0&1&6&4\end{array}\right],
\qquad P_{\mathcal{B}\leftarrow\mathcal{C}}=\begin{bmatrix}5&3\\6&4\end{bmatrix}.
\]
(بررسی: $5\vb_1+6\vb_2=(-7,9)=\vc_1$.) پس $P_{\mathcal{C}\leftarrow\mathcal{B}}=P_{\mathcal{B}\leftarrow\mathcal{C}}^{-1}=\frac12\begin{bmatrix}4&-3\\-6&5\end{bmatrix}$.
\end{solution}

\begin{example}
در فضای ماتریس‌های $2\times2$ مختلط، $P_{\mathcal{C}\leftarrow\mathcal{B}}$ را بیابید که
\[
\mathcal{B}=\left\{\begin{bmatrix}1&0\\0&1\end{bmatrix},\begin{bmatrix}0&1\\1&0\end{bmatrix},\begin{bmatrix}0&-i\\i&0\end{bmatrix},\begin{bmatrix}1&0\\0&-1\end{bmatrix}\right\},\quad
\mathcal{C}=\left\{\begin{bmatrix}1&0\\0&0\end{bmatrix},\begin{bmatrix}1&1\\0&0\end{bmatrix},\begin{bmatrix}1&1\\1&0\end{bmatrix},\begin{bmatrix}1&1\\1&1\end{bmatrix}\right\}.
\]
\end{example}
\begin{solution}
اگر $\begin{bmatrix}p&q\\r&s\end{bmatrix}=c_1\vc_1+\dots+c_4\vc_4$، از درایه‌های $(2,2)$، $(2,1)$، $(1,2)$ و $(1,1)$ به ترتیب $c_4=s$، $c_3=r-s$، $c_2=q-r$ و $c_1=p-q$. پس مختصات یک ماتریس نسبت به $\mathcal{C}$ برابر $(p-q,\ q-r,\ r-s,\ s)$ است و با اعمال آن روی چهار ماتریس $\mathcal{B}$:
\[
P_{\mathcal{C}\leftarrow\mathcal{B}}=\begin{bmatrix}1&-1&i&1\\0&0&-2i&0\\-1&1&i&1\\1&0&0&-1\end{bmatrix}.
\]
\end{solution}

\section{تبدیل خطی و تغییر پایه}

\begin{theorem}[ماتریس تبدیل در پایه‌ی دلخواه]\label{thm:matrix-of-map}
فرض کنید $T\in\mathcal{L}(V,W)$ و $\mathcal{E}=\{\ve_1,\dots,\ve_n\}$ و $\mathcal{B}=\{\vb_1,\dots,\vb_m\}$ پایه‌های $V$ و $W$ باشند. ماتریس
\[
[T]_{\mathcal{B}\leftarrow\mathcal{E}}=\begin{bmatrix}\coord{T\ve_1}{\mathcal{B}}&\cdots&\coord{T\ve_n}{\mathcal{B}}\end{bmatrix}\quad(m\times n)
\]
یکتا ماتریسی است که برای هر $\vu\in V$:
\begin{keyformula}
\[
\coord{T\vu}{\mathcal{B}}=[T]_{\mathcal{B}\leftarrow\mathcal{E}}\,\coord{\vu}{\mathcal{E}}.
\]
\end{keyformula}
\end{theorem}
\begin{proof}
اگر $\vu=\sum c_i\ve_i$، از خطی بودن $T$ و نگاشت مختصات: $\coord{T\vu}{\mathcal{B}}=\sum c_i\coord{T\ve_i}{\mathcal{B}}=[T]_{\mathcal{B}\leftarrow\mathcal{E}}\coord{\vu}{\mathcal{E}}$. یکتایی مثل قضیه‌ی \ref{thm:change-basis} است.
\end{proof}

پس ماتریس تغییر پایه همان ماتریس تبدیل همانی است: $P_{\mathcal{C}\leftarrow\mathcal{B}}=[I]_{\mathcal{C}\leftarrow\mathcal{B}}$.

\begin{example}
ماتریس مشتق‌گیری $T:P_3\to P_2$ را نسبت به پایه‌های $\mathcal{B}=\{x^3,x^2,x,1\}$ و $\mathcal{B}'=\{x^2,x,1\}$ بیابید.
\end{example}
\begin{solution}
مشتق پایه‌ها: $3x^2$، $2x$، $1$ و $0$ که مختصاتشان $(3,0,0)$، $(0,2,0)$، $(0,0,1)$ و $(0,0,0)$ است. پس
\[
[T]_{\mathcal{B}'\leftarrow\mathcal{B}}=\begin{bmatrix}3&0&0&0\\0&2&0&0\\0&0&1&0\end{bmatrix}.
\]
\end{solution}

\subsection*{تبدیل از فضا به خودش: $\mathcal{L}(V)$}

\begin{theorem}[تشابه]\label{thm:similarity}
فرض کنید $T:\R^n\to\R^n$ با $T(\vx)=A\vx$ و $\mathcal{B}=\{\vv_1,\dots,\vv_n\}$ پایه‌ای از $\R^n$ و $P=\begin{bmatrix}\vv_1&\cdots&\vv_n\end{bmatrix}$ باشد. آن‌گاه
\begin{keyformula}
\[
\coord{T\vx}{\mathcal{B}}=P^{-1}AP\,\coord{\vx}{\mathcal{B}},\qquad [T]_{\mathcal{B}}=P^{-1}[T]_{\mathcal{E}}P .
\]
\end{keyformula}
\end{theorem}
\begin{proof}
$\vx=P\coord{\vx}{\mathcal{B}}$ و $T\vx=A\vx=AP\coord{\vx}{\mathcal{B}}$. مختصات $T\vx$ نسبت به $\mathcal{B}$ برابر $P^{-1}T\vx=P^{-1}AP\coord{\vx}{\mathcal{B}}$ است.
\end{proof}

\begin{center}
\begin{tikzpicture}[>=Latex, box/.style={font=\large}]
  \node[box] (a) at (0,2) {$\R^n$}; \node[box] (b) at (4,2) {$\R^n$};
  \node[box] (c) at (0,0) {$\R^n$}; \node[box] (d) at (4,0) {$\R^n$};
  \draw[->, thick] (a)--(b) node[midway, above] {\small $[T]_{\mathcal{B}}$};
  \draw[->, thick] (c)--(d) node[midway, below] {\small $[T]_{\mathcal{E}}=A$};
  \draw[->, thick] (a)--(c) node[midway, left] {\small $P$};
  \draw[->, thick] (d)--(b) node[midway, right] {\small $P^{-1}$};
\end{tikzpicture}
\end{center}

ماتریس‌های $A$ و $P^{-1}AP$ را \term{متشابه}{Similar} می‌نامیم؛ هر دو یک تبدیل را در دو پایه نشان می‌دهند. با انتخاب پایه‌ی مناسب، $P^{-1}AP$ می‌تواند خیلی ساده‌تر از $A$ باشد.

\section{یکریختی و ماتریس}

مختصات‌گیری $\vx\mapsto\coord{\vx}{\mathcal{B}}$ یک یکریختی \enterm{Isomorphism} $V\to\F^n$ است (خطی طبق قضیه‌ی \ref{thm:coord-map}، و دوسویی \enterm{Bijective} چون نمایش یکتاست). از آن، یک روش استاندارد برای بررسی یکریخت \enterm{Isomorphic} بودن به دست می‌آید.

\begin{theorem}\label{thm:matrix-invertible-map}
اگر $\dim V=\dim W=n$ و $T\in\mathcal{L}(V,W)$، آن‌گاه $T$ وارون‌پذیر است اگر و تنها اگر $[T]_{\mathcal{B}\leftarrow\mathcal{E}}$ وارون‌پذیر باشد.
\end{theorem}
\begin{proof}
اگر $A=[T]_{\mathcal{B}\leftarrow\mathcal{E}}$ و $\varphi_{\mathcal{E}}$ و $\varphi_{\mathcal{B}}$ نگاشت‌های مختصات باشند، قضیه‌ی \ref{thm:matrix-of-map} می‌گوید $\varphi_{\mathcal{B}}\circ T=T_A\circ\varphi_{\mathcal{E}}$؛ یعنی $T=\varphi_{\mathcal{B}}^{-1}\circ T_A\circ\varphi_{\mathcal{E}}$. چون $\varphi$ها یکریختی‌اند، $T$ دوسویی است اگر و تنها اگر $T_A$ دوسویی باشد؛ یعنی $A$ وارون‌پذیر باشد (قضیه‌ی \ref{thm:imt-full}).
\end{proof}

\begin{example}
در مثال $e^x$ فصل ۱۱، ماتریس $T(a,b,c)=ae^x+bxe^x+cx^2e^x$ نسبت به پایه‌ی استاندارد و پایه‌ی $\{e^x,xe^x,x^2e^x\}$ برابر $I_3$ است؛ پس وارون‌پذیر است و $T$ یکریختی است.
\end{example}

\section{جمع‌بندی و تمرین}

\begin{summary}
\begin{itemize}
  \item $P_{\mathcal{C}\leftarrow\mathcal{B}}$ ستون‌هایش $\coord{\vb_j}{\mathcal{C}}$ است و $\coord{\vx}{\mathcal{C}}=P_{\mathcal{C}\leftarrow\mathcal{B}}\coord{\vx}{\mathcal{B}}$؛ وارون‌پذیر است و $P_{\mathcal{B}\leftarrow\mathcal{C}}=P_{\mathcal{C}\leftarrow\mathcal{B}}^{-1}$. در $\R^n$: $P_{\mathcal{C}\leftarrow\mathcal{B}}=P_{\mathcal{C}}^{-1}P_{\mathcal{B}}$.
  \item $[T]_{\mathcal{B}\leftarrow\mathcal{E}}$ ستون‌هایش $\coord{T\ve_j}{\mathcal{B}}$ است و $\coord{T\vu}{\mathcal{B}}=[T]_{\mathcal{B}\leftarrow\mathcal{E}}\coord{\vu}{\mathcal{E}}$.
  \item برای $T:\R^n\to\R^n$: $[T]_{\mathcal{B}}=P^{-1}AP$؛ ماتریس‌های متشابه یک تبدیل را در دو پایه نشان می‌دهند.
  \item $T$ وارون‌پذیر است $\iff$ ماتریسش (در هر جفت پایه) وارون‌پذیر است.
\end{itemize}
\end{summary}

\subsection*{تمرین‌ها}
\begin{enumerate}
  \item برای $\mathcal{B}=\{(1,1),(1,-1)\}$ و پایه‌ی استاندارد $\mathcal{E}$ در $\R^2$، $P_{\mathcal{E}\leftarrow\mathcal{B}}$ و $\coord{(3,1)}{\mathcal{B}}$ را بیابید.
  \item ماتریس تبدیل $T(x,y)=(2x,\,3y)$ را در پایه‌ی $\mathcal{B}=\{(1,1),(1,-1)\}$ بیابید.
  \item ماتریس مشتق‌گیری $P_2\to P_2$ را نسبت به پایه‌ی $\{1,x,x^2\}$ بیابید.
\end{enumerate}

\subsection*{پاسخ‌های کوتاه}
\begin{enumerate}
  \item $P_{\mathcal{E}\leftarrow\mathcal{B}}=\begin{bmatrix}1&1\\1&-1\end{bmatrix}$ و $\coord{(3,1)}{\mathcal{B}}=P^{-1}(3,1)=(2,1)$.
  \item $P^{-1}AP$ با $A=\mathrm{diag}(2,3)$ و $P=\begin{bmatrix}1&1\\1&-1\end{bmatrix}$ برابر $\frac12\begin{bmatrix}5&-1\\-1&5\end{bmatrix}$ است.
  \item $T(1)=0$، $T(x)=1$، $T(x^2)=2x$؛ ماتریس $\begin{bmatrix}0&1&0\\0&0&2\\0&0&0\end{bmatrix}$.
\end{enumerate}
